\documentclass[11pt,a4paper]{article}
\usepackage[utf8]{inputenc}
\usepackage[T1]{fontenc}
\usepackage[ngerman]{babel}
\usepackage{geometry}
\usepackage{hyperref}
\geometry{margin=2.5cm}
\setlength{\parindent}{0pt}
\setlength{\parskip}{0.7em}
\begin{document}
\begin{center}
{\Large KI-Bericht -- Web-Technologien, Übungsblatt 1}\\
Fabio Schmickl\\
621.901 Web-Technologien UE, WS 2026
\end{center}

\section*{Anbieter, Modell und Version}
Anbieter laut \texttt{hermes status}: ChatGPT or Codex Subscription.\\
Modell laut \texttt{hermes status}: \texttt{gpt-5.5}.\\
Eine separate Modellversionsnummer wurde von \texttt{hermes status} nicht ausgegeben.

\section*{Reasoning-Effort}
High.

\section*{Tools und Plug-ins}
Verwendet wurden Hermes Agent, die Skills \texttt{moodle}, \texttt{solve-and-explain}, \texttt{github-operations}, \texttt{pdf} und \texttt{software-engineering-workflows}, ausserdem Terminal-Werkzeuge wie \texttt{git}, \texttt{gh}, \texttt{pdftotext}, \texttt{pdftoppm}, \texttt{curl}, \texttt{python3} und \texttt{pdflatex}.

\section*{Kontext}
Verwendeter Kontext: die Moodle-Angabe \texttt{Assignment01.pdf}, der Moodle-Kontext zu Übungsblatt 1, die Zusatzdatei \texttt{files.zip}, die KI-Erlaubnisdatei \texttt{62190126W-ki-erlaubnis.md} und das Kurs-Repository \texttt{Fafoooo/aau\_26W\_Webtechnologien}.

\section*{Tatsächlicher Ablauf}
Die Angabe wurde aus Moodle übernommen, als Text extrahiert und zusätzlich als PDF-Seitenbild geprüft, damit die Layout-Vorlagen für Exercise 3, 8 und 9 sichtbar waren. Danach wurden die Zusatzdateien aus \texttt{files.zip} entpackt. Für die HTTP-Aufgabe wurden aktuelle Header- und Ressourcenmessungen für AAU und ORF durchgeführt. Anschließend wurden die HTML-, CSS- und Antwortdateien pro Aufgabe in eigenen Ordnern erstellt. Die Dateien wurden mit einem lokalen Prüfsktipt auf vorhandene Pflichtdateien, Links, Formularattribute, Tabellenstruktur, Galerie-Meta-Refresh und CSS-Anforderungen geprüft. Danach wurde diese Erklärung und dieser KI-Bericht erzeugt.

\section*{Prompt-Verlauf sinngemäß}
Der automatisierte Moodle-Solver-Auftrag lautete sinngemäß: Repo holen, Angabe und Zusatzdateien ablegen, Übungsblatt 1 für Web-Technologien lösen, für jede Teilaufgabe lauffähigen Code und Antworten erzeugen, eine Erklärung für Fabio schreiben, einen KI-Bericht mit echten Angaben erstellen, testen, committen, pushen und den Moodle-Solver-Status aktualisieren. Tatsächliche Teilaufträge an die KI waren: Angabe erfassen, Zusatzdateien aus Moodle holen, HTTP-Daten messen, HTML/CSS/Formular-/Tabellenlösungen erstellen, CSS-Selektoren erklären, Dateien prüfen, unabhängige Prüfung durchführen und Bericht erzeugen.

\section*{Eigene Bearbeitung}
Dieser Abschnitt ist von Fabio auszufüllen. Hier kann dokumentiert werden, was selbst nachvollzogen, getestet, geändert oder ergänzt wurde.

\vspace{3cm}
\hrule
\vspace{0.4cm}
Notizen eigene Bearbeitung:

\vspace{4cm}
\hrule

\end{document}
